%%%%%%%%%%%%%%%%%%%%%%%%%%%%%%%%%%%%%%%%%%%%%%%%%%%%%%%%%%%%%%%%%%%%%%%%%%%%%%%%
% ACF-DOC-039
% Multi-Hazard Decision Support & Emergency Management
% Chapter 01 - Multi-Hazard Architecture
%%%%%%%%%%%%%%%%%%%%%%%%%%%%%%%%%%%%%%%%%%%%%%%%%%%%%%%%%%%%%%%%%%%%%%%%%%%%%%%%

\section{Multi-Hazard Architecture}

\subsection{Executive Summary}

The Atmospheric Complexity Framework (ACF) introduces a global
Multi-Hazard Decision Support and Emergency Management System designed
to transform Earth system observations and numerical predictions into
operational intelligence.

The objective of this system is not only to forecast atmospheric,
hydrological, oceanic and environmental hazards, but also to evaluate
their potential impacts and support operational decision making through
the Atmospheric Weather Center Interface (AWCI).

The ACF Multi-Hazard Engine integrates:

\begin{itemize}

\item Numerical Weather Prediction systems

\item Earth observation networks

\item Data assimilation technologies

\item Artificial intelligence models

\item Earth System Digital Twin

\item Risk assessment algorithms

\item Operational alert management

\end{itemize}

The final objective is the creation of a unified global platform capable
of monitoring, predicting and managing multiple interacting hazards.

%%%%%%%%%%%%%%%%%%%%%%%%%%%%%%%%%%%%%%%%%%%%%%%%%%%%%%%%%%%%%%%%%%%%%%%%%%%%%%%%

\subsection{Introduction}

Natural hazards are complex Earth system phenomena resulting from
interactions between atmosphere, ocean, land surface, cryosphere and
human activities.

Traditional warning systems often analyse individual hazards
independently. However, modern Earth system prediction requires a
multi-hazard approach because several hazards can occur simultaneously.

Examples include:

\begin{itemize}

\item Tropical cyclones producing extreme rainfall and flooding

\item Heatwaves combined with drought and wildfire risk

\item Severe storms producing tornadoes, hail and flash floods

\item Dust storms affecting air quality and transportation

\item Solar storms affecting technological infrastructure

\end{itemize}

The ACF architecture therefore treats hazards as interconnected
components of a global Earth system.

%%%%%%%%%%%%%%%%%%%%%%%%%%%%%%%%%%%%%%%%%%%%%%%%%%%%%%%%%%%%%%%%%%%%%%%%%%%%%%%%

\subsection{ACF Multi-Hazard Intelligence Concept}

The ACF Multi-Hazard Intelligence Engine follows a complete information
chain:

\[
Observation
\rightarrow
Assimilation
\rightarrow
Digital Twin
\rightarrow
Hazard Detection
\rightarrow
Risk Analysis
\rightarrow
Decision Support
\rightarrow
Alert
\]

Each stage continuously updates the representation of the Earth system.

%%%%%%%%%%%%%%%%%%%%%%%%%%%%%%%%%%%%%%%%%%%%%%%%%%%%%%%%%%%%%%%%%%%%%%%%%%%%%%%%

\subsection{Global Hazard Classification}

The ACF hazard framework includes several categories.

\subsubsection{Atmospheric Hazards}

\begin{itemize}

\item Severe thunderstorms

\item Tornadoes

\item Lightning

\item Hail

\item Derechos

\item Heavy precipitation

\item Wind storms

\item Fog

\item Dust storms

\end{itemize}


\subsubsection{Hydrological Hazards}

\begin{itemize}

\item River floods

\item Flash floods

\item Urban flooding

\item Coastal flooding

\item Soil saturation events

\end{itemize}


\subsubsection{Climate Hazards}

\begin{itemize}

\item Heatwaves

\item Cold waves

\item Drought

\item Extreme climate anomalies

\end{itemize}


\subsubsection{Ocean Hazards}

\begin{itemize}

\item Tropical cyclones

\item Storm surge

\item Extreme waves

\item Marine hazards

\end{itemize}


\subsubsection{Environmental Hazards}

\begin{itemize}

\item Wildfires

\item Air pollution

\item Aerosols

\item Chemical hazards

\end{itemize}


\subsubsection{Space Weather Hazards}

\begin{itemize}

\item Solar flares

\item Geomagnetic storms

\item Radiation events

\item Ionospheric disturbances

\end{itemize}

%%%%%%%%%%%%%%%%%%%%%%%%%%%%%%%%%%%%%%%%%%%%%%%%%%%%%%%%%%%%%%%%%%%%%%%%%%%%%%%%

\subsection{Multi-Hazard System Architecture}

The global architecture is represented as:

\begin{center}

\begin{tikzpicture}

\node (obs) [draw, rounded corners] 
{Global Observation Network};

\node (assim) [draw, rounded corners, below=1cm of obs]
{Data Assimilation System};

\node (twin) [draw, rounded corners, below=1cm of assim]
{Earth System Digital Twin};

\node (hazard) [draw, rounded corners, below=1cm of twin]
{Multi-Hazard Detection Engine};

\node (risk) [draw, rounded corners, below=1cm of hazard]
{AI Risk Assessment Engine};

\node (decision) [draw, rounded corners, below=1cm of risk]
{Decision Support System};

\node (awci) [draw, rounded corners, below=1cm of decision]
{AWCI Operational Center};


\draw[->] (obs) -- (assim);
\draw[->] (assim) -- (twin);
\draw[->] (twin) -- (hazard);
\draw[->] (hazard) -- (risk);
\draw[->] (risk) -- (decision);
\draw[->] (decision) -- (awci);

\end{tikzpicture}

\end{center}

%%%%%%%%%%%%%%%%%%%%%%%%%%%%%%%%%%%%%%%%%%%%%%%%%%%%%%%%%%%%%%%%%%%%%%%%%%%%%%%%

\subsection{Earth System Digital Twin Coupling}

The Multi-Hazard Engine receives continuously updated information from
the ACF Digital Twin.

The digital representation includes:

\begin{itemize}

\item Atmospheric state

\item Ocean conditions

\item Land surface parameters

\item Soil moisture

\item Snow and ice conditions

\item Chemical composition

\item Human exposure information

\end{itemize}

This enables prediction of compound hazards and cascading impacts.

%%%%%%%%%%%%%%%%%%%%%%%%%%%%%%%%%%%%%%%%%%%%%%%%%%%%%%%%%%%%%%%%%%%%%%%%%%%%%%%%

\subsection{Risk Assessment Framework}

The ACF risk model is based on:

\[
Risk =
Hazard
\times
Exposure
\times
Vulnerability
\times
Confidence
\]

where:

\begin{itemize}

\item Hazard represents the physical intensity

\item Exposure represents affected elements

\item Vulnerability represents sensitivity

\item Confidence represents forecast reliability

\end{itemize}

%%%%%%%%%%%%%%%%%%%%%%%%%%%%%%%%%%%%%%%%%%%%%%%%%%%%%%%%%%%%%%%%%%%%%%%%%%%%%%%%
%%%%%%%%%%%%%%%%%%%%%%%%%%%%%%%%%%%%%%%%%%%%%%%%%%%%%%%%%%%%%%%%%%%%%%%%%%%%%%%%
\subsection{Artificial Intelligence Hazard Intelligence Engine}

The ACF Multi-Hazard system integrates artificial intelligence methods
to improve hazard detection, prediction and impact assessment.

The AI Hazard Intelligence Engine combines:

\begin{itemize}

\item Machine Learning classification

\item Deep Neural Networks

\item Convolutional Neural Networks for spatial analysis

\item Transformer models for temporal prediction

\item Graph Neural Networks for Earth system interactions

\item Physics-Informed Neural Networks

\end{itemize}

The AI system does not replace physical models. It operates as a hybrid
intelligence layer combining numerical simulation and data-driven
learning.

The general hybrid prediction concept is:

\[
Forecast_{ACF}
=
Physics_{Model}
+
AI_{Correction}
+
Observation_{Assimilation}
\]

%%%%%%%%%%%%%%%%%%%%%%%%%%%%%%%%%%%%%%%%%%%%%%%%%%%%%%%%%%%%%%%%%%%%%%%%%%%%%%%%

\subsection{Hazard Detection Pipeline}

The operational detection chain contains several processing stages.

\begin{enumerate}

\item Data acquisition

\item Quality control

\item Data assimilation

\item Feature extraction

\item Hazard identification

\item Probability estimation

\item Risk calculation

\item Alert generation

\end{enumerate}


The pipeline receives information from:

\begin{itemize}

\item Weather satellites

\item Radar networks

\item Surface stations

\item Radiosondes

\item Aircraft observations

\item Ocean platforms

\item Numerical models

\item Internet of Things sensors

\end{itemize}

%%%%%%%%%%%%%%%%%%%%%%%%%%%%%%%%%%%%%%%%%%%%%%%%%%%%%%%%%%%%%%%%%%%%%%%%%%%%%%%%

\subsection{Multi-Hazard Risk Matrix}

ACF uses a multi-dimensional risk matrix.

\begin{center}

\begin{tabular}{llll}

\toprule

Parameter & Description & Range & Unit \\

\midrule

Hazard intensity & Physical magnitude & 0-5 & Index \\

Probability & Event likelihood & 0-100 & \% \\

Exposure & Population/assets affected & 0-100 & \% \\

Vulnerability & Sensitivity factor & 0-100 & \% \\

Confidence & Forecast reliability & 0-100 & \% \\

\bottomrule

\end{tabular}

\end{center}


The final operational risk score is calculated as:

\[
R_s =
I_h
\times
P_e
\times
V_u
\times
C_f
\]

where:

\begin{itemize}

\item $I_h$ = hazard intensity

\item $P_e$ = probability of occurrence

\item $V_u$ = vulnerability coefficient

\item $C_f$ = forecast confidence

\end{itemize}

%%%%%%%%%%%%%%%%%%%%%%%%%%%%%%%%%%%%%%%%%%%%%%%%%%%%%%%%%%%%%%%%%%%%%%%%%%%%%%%%

\subsection{Operational Decision Support Framework}

The ACF Decision Support System transforms scientific information
into operational recommendations.

The system provides:

\begin{itemize}

\item Hazard monitoring

\item Risk ranking

\item Alert recommendation

\item Impact estimation

\item Response prioritization

\item Situation awareness

\end{itemize}


The operational workflow is:

\[
Prediction
\rightarrow
Risk Evaluation
\rightarrow
Decision
\rightarrow
Action
\]

%%%%%%%%%%%%%%%%%%%%%%%%%%%%%%%%%%%%%%%%%%%%%%%%%%%%%%%%%%%%%%%%%%%%%%%%%%%%%%%%

\subsection{AWCI Operational Integration}

The Atmospheric Weather Center Interface (AWCI) represents the final
operational layer of the ACF Multi-Hazard system.

AWCI provides:

\begin{itemize}

\item Global hazard maps

\item Real-time monitoring panels

\item Warning levels

\item Forecast confidence indicators

\item Event tracking

\item Emergency communication tools

\end{itemize}


The AWCI dashboard receives:

\begin{itemize}

\item Hazard probabilities

\item Geographic risk areas

\item Time evolution

\item Model uncertainty

\item AI explanations

\end{itemize}

%%%%%%%%%%%%%%%%%%%%%%%%%%%%%%%%%%%%%%%%%%%%%%%%%%%%%%%%%%%%%%%%%%%%%%%%%%%%%%%%

\subsection{Operational Alert Levels}

The ACF warning system defines several operational levels.

\begin{center}

\begin{tabular}{lll}

\toprule

Level & Meaning & Action \\

\midrule

Green & Normal situation & Monitoring \\

Yellow & Potential hazard & Preparation \\

Orange & Significant hazard & Operational response \\

Red & Extreme hazard & Emergency action \\

Purple & Catastrophic event & Maximum response \\

\bottomrule

\end{tabular}

\end{center}

%%%%%%%%%%%%%%%%%%%%%%%%%%%%%%%%%%%%%%%%%%%%%%%%%%%%%%%%%%%%%%%%%%%%%%%%%%%%%%%%

\subsection{Real-Time Processing Architecture}

The operational system requires continuous processing.

\begin{tikzpicture}

\node (data) [draw, rounded corners]
{Real-Time Data Streams};

\node (hpc) [draw, rounded corners, below=1cm of data]
{HPC Processing Layer};

\node (ai) [draw, rounded corners, below=1cm of hpc]
{AI Analysis};

\node (risk) [draw, rounded corners, below=1cm of ai]
{Risk Engine};

\node (alert) [draw, rounded corners, below=1cm of risk]
{Alert System};


\draw[->] (data) -- (hpc);
\draw[->] (hpc) -- (ai);
\draw[->] (ai) -- (risk);
\draw[->] (risk) -- (alert);

\end{tikzpicture}

%%%%%%%%%%%%%%%%%%%%%%%%%%%%%%%%%%%%%%%%%%%%%%%%%%%%%%%%%%%%%%%%%%%%%%%%%%%%%%%%

\subsection{Future Development}

Future versions of ACF will introduce:

\begin{itemize}

\item Global probabilistic hazard forecasting

\item Autonomous AI emergency assistant

\item High-resolution urban risk prediction

\item Climate change impact assessment

\item Digital Twin based scenario simulation

\item Fully automated decision support

\end{itemize}

%%%%%%%%%%%%%%%%%%%%%%%%%%%%%%%%%%%%%%%%%%%%%%%%%%%%%%%%%%%%%%%%%%%%%%%%%%%%%%%%
